\section{Related Work}
\label{sec: related work}
% \subsection{Dynamic Graph Systems}
% \label{subsec:dynamic}
% ============================================================
% \YJ{Discuss other related works: CuckooGraph, etc.}

\textit{Dynamic graph computation.}
% Many systems have been proposed to support efficient storage-level
% updates on evolving graphs.
% We have extensively discussed Terrace~\cite{pandey2021terrace}, 
% Spruce~\cite{shi2024spruce}, and Aspen~\cite{dhulipala2019low}
% throughout the paper. 
Besides the extensively discussed Terrace~\cite{pandey2021terrace},
Spruce~\cite{shi2024spruce}, Aspen~\cite{dhulipala2019low},
GraphBolt~\cite{mariappan2019graphbolt}, and DZiG~\cite{mariappan2021dzig},
Stinger~\cite{ediger2012stinger} aggregates edge lists into
fixed-size blocks with a hash-table index for point queries.
GraphOne~\cite{kumar2020graphone} combines an adjacency list with
an edge log that is periodically merged in batches.
Sortledton~\cite{fuchs2023sortledton} replaces edge lists with an
unrolled skip list offering $O(\log|E|)$ update and lookup with
transactional support.
Teseo~\cite{de2021teseo} uses packed memory arrays as leaves of
B+ trees with optimistic locking for concurrent access.
CuckooGraph~\cite{fan2025cuckoograph} proposes a hash-array-based
structure that adaptively resizes through transformation and denylist
techniques.
Other works include LiveGraph~\cite{zhu2019livegraph} and
LLAMA~\cite{macko2015llama}.



\textit{Query-dependent incremental computation.}
A separate line of work attempts to reuse prior computational results
across updates~\cite{besta2021practice,ramalingam1996computational}.
\textit{Khanda et al.}~\cite{khanda2021parallel} maintain an SSSP
tree rooted at a fixed source after each update, recomputing only the
affected subtree.
\textit{Kickstarter}~\cite{vora2017kickstarter} retains approximate monotonic
results and propagates corrections for invalidated values.
% \cq{ 
\textit{GraphBolt}~\cite{mariappan2019graphbolt} uses dependency-driven
synchronous processing: it tracks per-vertex aggregation dependencies across
BSP iterations and propagates only the delta caused by mutations, supporting
algorithms including PageRank, belief propagation, SSSP, and BFS.

\textit{DZiG}~\cite{mariappan2021dzig} extends GraphBolt with sparsity-aware
processing that skips zero-contribution edges during incremental propagation,
reducing redundant work on sparse updates.
\textit{RisGraph}~\cite{feng2021risgraph} targets sub-millisecond per-update
analysis by streaming graph mutations into a compact active-vertex frontier.
\textit{GraphFly}~\cite{chen2022graphfly} uses asynchronous dependency-flow
execution to reduce synchronization overhead in streaming graph processing.
\textit{IncBoost}~\cite{yin2024incboost} focuses on incremental processing
under edge deletions and weight updates.
\textit{Fan et al.}~\cite{fan2021incrementalizing} propose a general framework
for incrementalizing arbitrary graph algorithms by detecting affected
regions across updates.
\textit{Bauer and Wagner}~\cite{bauer2009batch} study batch dynamic SSSP
experimentally, showing reuse is only effective when the affected
subgraph is small.

For betweenness centrality specifically, \textit{Green et al.}~\cite{green2012fast}
track path-count changes caused by each edge insertion or deletion to maintain
an incremental BC score, while \textit{Nasre et al.}~\cite{Nasre2014IncBC}
give an asymptotically faster incremental BC algorithm that exploits DAG structure
to localize the repair.
All of these systems maintain algorithm-specific state (e.g., per-vertex
aggregation values, SSSP trees, path-count tables, or monotonic approximations)
that must be reset or rebuilt when the algorithm or source changes.
GIR instead maintains \emph{algorithm-agnostic} structural invariants
(WCC, SCC, BiCC) that benefit BFS, SSSP, and BC simultaneously, with no
per-algorithm reconfiguration.
% }

\textit{Theoretical dynamic graph algorithms.}
A substantial body of theory addresses fully dynamic connectivity; we refer
to Hanauer et al.~\cite{Hanauer2022DynSurvey} for a comprehensive survey.
\textit{Holm et al.}~\cite{holm2001poly} give poly-logarithmic deterministic
algorithms for dynamic connectivity, MST, 2-edge connectivity, and
biconnectivity.
\textit{Wulff-Nilsen}~\cite{wulff2013faster} improves the deterministic
connectivity bound.
\textit{De Man et al.}~\cite{de2025towards} study batch-dynamic
connectivity in a practical parallel setting, targeting throughput on
large batches of concurrent edge updates.
These algorithms target worst-case update complexity under arbitrary adversarial
sequences.
GIR's maintenance is not designed to match these worst-case guarantees; instead,
it exploits the SPU five-case classification to achieve \emph{low average-case
cost} on real-world graphs, where most updates fall into cheap cases (Case 1
or Case 4) and full structural repair is rare.
Unlike purely theoretical constructions, GIR integrates maintenance with a
two-level parallel execution framework and is evaluated end-to-end on real
graph workloads.
% These approaches can be highly efficient when the same algorithm is
% repeated from the same source, but they are tightly coupled to a
% specific algorithm and, in some cases, a specific source node.
% In practice, dynamic graph workloads involve ad hoc queries with
% varying sources, targets, and algorithms~\cite{Hanauer2022DynSurvey},
% so such query-dependent systems offer limited reuse.


% \textbf{Intermediate representations and graph frameworks.}
\textit{Intermediate representations (IRs)}, such as LLVM
IR~\cite{Lattner2004LLVM}, decouple source languages from back-end
optimization passes, enabling reusable analysis across many
compilation stages.
% GIR is directly inspired by this principle: just as a compiler IR is
% computed once and reused across many passes, GIR maintains structural
% invariants once and makes them available to many downstream
% algorithms.
Several graph domain-specific languages adopt a related
algorithm-execution separation.
GraphIt~\cite{Zhang2018GraphIt} separates algorithm descriptions from
scheduling decisions.
Galois~\cite{nguyen2013lightweight} decouples
operator logic from parallel execution policies.
Ligra~\cite{Shun2013Ligra}, Green-Marl~\cite{hong2012green},
Gemini~\cite{zhu2016gemini}, and Pregel~\cite{malewicz2010pregel}
target shared- and distributed-memory execution on static snapshots.

% Unlike these frameworks, GIR is not a programming abstraction or a
% scheduling layer.
% Instead, it is a persistent structural layer maintained
% incrementally under continuous updates and exposed directly to
% downstream algorithms.
% No existing system simultaneously maintains multi-granular,
% query-agnostic structural invariants, repairs them incrementally
% under dynamic updates, and exposes them for task-parallel query
% execution. This is the gap GIR fills.


% Terrace~\cite{pandey2021terrace},
% Aspen~\cite{dhulipala2019low}, and Spruce~\cite{shi2024spruce}
% (evaluated as baselines in Section~\ref{sec:experiment}) take
% complementary storage-layer approaches.
% All of these systems optimize for update throughput at the storage
% layer but provide no persistent structural invariants for downstream
% analytics.

% CuckooGraph~\cite{fan2025cuckoograph} proposes a hash-array-based structure for large-scale dynamic graphs that adaptively resizes through TRANSFORMATION and DENYLIST techniques, achieving high insertion throughput with low memory overhead.
% Like Terrace, Spruce, and Aspen, CuckooGraph is optimized purely for storage-level update performance and does not maintain any structural invariants for downstream analytics.
% Several other adjacency-list-based systems have been proposed for dynamic graph storage.
% Stinger~\cite{ediger2012stinger} aggregates edge lists into fixed-size blocks and uses a
% hash table as a secondary index to improve point query efficiency.
% GraphOne~\cite{kumar2020graphone} combines an adjacency list with an edge log array:
% updates are appended to the log and periodically merged into the adjacency list in batches,
% supporting real-time analytics on evolving graphs.
% LiveGraph~\cite{zhu2019livegraph} stores edges in purely sequential transactional edge logs,
% optimizing sequential scan performance for read-heavy workloads.
% Sortledton~\cite{fuchs2023sortledton} replaces edge lists with an unrolled skip list,
% providing $O(\log|E|)$ lookup, insertion, and deletion with transactional support.
% Teseo~\cite{de2021teseo} uses packed memory arrays as leaves of B$^+$ trees and
% supports concurrent reads and writes through optimistic locking.
% LLAMA~\cite{macko2015llama} takes a snapshot-oriented approach,
% storing graph versions as large multiversioned arrays to support
% efficient analytics on historical snapshots.
% Like CuckooGraph, Terrace, Spruce, and Aspen, all of these systems target update throughput
% at the storage layer and provide no persistent structural invariants for downstream analytics.

% Dynamic graph systems must support two operations: integrating
% structural updates quickly and answering downstream queries
% efficiently.
% Existing systems focus primarily on the first operation.
% We survey the most relevant systems and explain their limitations
% for the second operation.

% \textbf{Graph storage systems.} 
% % \YJ{They should not be called "storage system".}
% Most dynamic graph systems propose specialized in-memory data
% structures that balance update cost, query efficiency,
% and memory usage.
% Static representations such as CSR and CSC~\cite{wheatman2018packed}
% store adjacency information in contiguous arrays and enable
% efficient traversal, but edge insertions and deletions require
% rebuilding large arrays, which is prohibitively expensive
% under frequent updates.

% Terrace~\cite{pandey2021terrace} targets skewed streaming graphs.
% It introduces a hierarchical adjacency representation in which
% a vertex's neighbors are stored using different data structures
% depending on its degree.
% As degrees evolve, the system dynamically reorganizes storage
% to improve locality and update efficiency on power-law graphs.
% However, Terrace focuses entirely on storage and does not
% maintain any structural invariants for downstream computation.

% Spruce~\cite{shi2024spruce} combines compact indexing with
% update-friendly adjacency organization.
% It integrates dynamic buffers with CSR-like sorted blocks and
% introduces a read-optimized concurrency protocol that supports
% efficient simultaneous reads and writes.
% Like Terrace, Spruce targets update throughput and leaves
% downstream analytics to be rerun from scratch after each update.

% Aspen~\cite{dhulipala2019low} represents graphs using compressed purely
% functional trees and exposes immutable snapshot views for
% analytics.
% Each update produces a new graph version, enabling concurrent
% readers on consistent snapshots while preserving structural
% sharing across versions.
% This versioning design is well suited for analytics on static
% snapshots, but it does not maintain structural invariants
% incrementally and requires full recomputation for downstream
% algorithms after each update.

% \YJ{Avoid duplicates. If they have been discussed in detail earlier, we can simply mention them here.}

% \textbf{Query-dependent computation systems.}
% Another line of work attempts to reuse prior computational
% results directly~\cite{besta2021practice,ramalingam1996computational}.
% For example, \textit{Khanda et al.}~\cite{khanda2021parallel} maintains an SSSP tree
% rooted at a fixed source after each query. %people--> according to the first
% After a graph update, it identifies the affected subtree and
% recomputes only that portion.
% Kickstarter~\cite{vora2017kickstarter} retains approximate monotonic
% computation results across updates, trimming values that have been
% invalidated by edge deletions and propagating corrections only
% within the affected region.
% Green et al.~\cite{green2012fast} maintain an incremental
% betweenness centrality score by tracking only the path-count
% changes caused by each edge insertion or deletion.
% Fan et al.~\cite{fan2021incrementalizing} propose a general framework
% for incrementalizing arbitrary graph algorithms by detecting affected
% regions and reusing unaffected partial results across updates.
% Bauer and Wagner~\cite{bauer2009batch} study batch dynamic SSSP algorithms
% experimentally, showing that incremental updates outperform recomputation
% only when the affected subgraph remains small relative to the full graph.
% These approaches are efficient when the same query is repeated
% under the same algorithm.
% However, each is tightly coupled to a specific algorithm and,
% in some cases, a specific source node.
% If the query algorithm or source changes, the cached state
% provides no benefit.
% In practice, dynamic graph workloads involve ad hoc queries
% with varying sources, targets, and algorithms~\cite{Hanauer2022DynSurvey},
% so such query-dependent systems offer limited reuse.

% % \YJ{Double check which 2-extreme does each method belong to.}

% % \textbf{GIR's position.}
% % GIR differs from all of these systems in one key respect.
% % Storage systems such as Terrace, Spruce, and Aspen are
% % history-blind: they discard all prior structural computation
% % after each update.
% % Query-dependent systems preserve results but tie them to
% % a specific query and source.
% % GIR occupies the middle ground: it maintains a small set of
% % query-agnostic structural invariants that persist across updates
% % and benefit many downstream algorithms without being tied to
% % any specific query or source.
% % GIR is designed to operate on top of existing storage systems
% % rather than replace them.
% % It can be integrated with formats such as Terrace, Aspen,
% % Spruce, and others~\cite{bender2007adaptive,king2016dynamic,islam2023dgap} as a
% % structural layer above the storage layer.

% % % ============================================================
% % \subsection{Intermediate Representations for Graphs}
% % ============================================================

% % \textbf{Intermediate representation in compilers.}
% % Compiler intermediate representations (IRs), such as LLVM IR~\cite{Lattner2004LLVM}, decouple source languages from optimization passes.
% % This abstraction enables reusable analysis and transformation across different compilation stages.\YJ{Discuss any works that either call or relate to IR for Graph} Several graph DSLs and compilers, such as GraphIt~\cite{Zhang2018GraphIt}, provide scheduling abstractions that separate algorithm specification from execution optimization. \YJ{Only one? Any others?}
% % However, these systems primarily focus on optimizing computation rather than maintaining persistent structural invariants under dynamic updates.

% % ============================================================

% % \textbf{Inter}
% \textbf{Intermediate representations (IRs)}, 
% % in compilers.}
% % Compiler intermediate representations (IRs), 
% such as LLVM
% IR~\cite{Lattner2004LLVM}, decouple source languages from
% back-end optimization passes.
% This abstraction enables reusable analysis and transformation
% across compilation stages.
% A key property of compiler IRs is that they are computed once
% and reused across many passes, rather than recomputed from
% source on every pass.
% % \textbf{Graph DSLs and scheduling abstractions.}
% Several graph domain-specific languages adopt a similar
% separation between algorithm specification and execution
% strategy.
% GraphIt~\cite{Zhang2018GraphIt} separates graph algorithm
% descriptions from scheduling decisions such as direction
% optimization, NUMA-awareness, and cache tiling.
% Galois~\cite{nguyen2013lightweight} provides a worklist-based runtime
% that separates operator logic from parallel execution policies.
% Ligra~\cite{Shun2013Ligra} provides a unified interface for
% vertex-subset-based graph algorithms on shared memory.
% Green-Marl~\cite{hong2012green} provides a domain-specific language
% with a property graph model that compiles to efficient parallel C++ code.
% Gemini~\cite{zhu2016gemini} targets distributed-memory execution with
% chunk-based graph partitioning.
% Pregel~\cite{malewicz2010pregel} introduces a vertex-centric bulk synchronous parallel model for large-scale graph processing.
% These systems improve programmability and performance for
% static graph analytics.
% However, they operate on fixed graph snapshots~\cite{mccune2015thinking}.
% They do not maintain persistent structural invariants across
% updates, and they require each downstream algorithm to
% rebuild structural knowledge from scratch.

% \textbf{GIR's distinction.}
% GIR adopts the IR philosophy and applies it to dynamic graphs.
% Like a compiler IR, GIR is computed incrementally and reused
% across many downstream operations.
% Unlike graph DSLs, GIR is not a programming abstraction.
% It is a persistent structural layer that is maintained under
% continuous updates and that exposes structural invariants
% directly to downstream algorithms.
% No prior dynamic graph system maintains a multi-granular,
% query-agnostic structural layer of this kind.
% This is the gap that GIR fills.